% CVPR 2027 submission; template: https://github.com/cvpr-org/author-kit (CVPR2027-v1 kit)
\documentclass[10pt,twocolumn,letterpaper]{article}

%%%%%%%%% PAPER TYPE
% \usepackage{cvpr}              % camera-ready
\usepackage[review]{cvpr}      % review version
% \usepackage[pagenumbers]{cvpr} % arXiv version

\usepackage{booktabs}
\usepackage{array}
\usepackage{graphicx}
\usepackage{multirow}
\usepackage[table]{xcolor}
\usepackage{amsmath}
\usepackage{url}
\usepackage{tikz}
\usetikzlibrary{positioning,arrows.meta,patterns}
\usepackage{microtype}  % kit preamble.tex: tighter layout
\definecolor{ourteal}{HTML}{0D5C63}
\definecolor{memgreen}{HTML}{2F8F83}
\definecolor{lexorange}{HTML}{C4841D}

% Submission-ready: debug macros neutralised.
\newcommand{\vnum}[1]{#1}
\newcommand{\todo}[1]{}

\definecolor{cvprblue}{rgb}{0.21,0.49,0.74}
%% pagebackref (recommended by cvpr.sty for the review copy) costs one extra pdflatex pass
%% (4 instead of 3). It is enabled only where refs.bib exists (local builds, which produce the
%% submission PDF) and skipped on Overleaf (free plan, 10 s compile limit).
%% It needs no BibTeX: the bibliography itself always comes from refs_precompiled.bbl when
%% that file exists (see the switch before \end{document}). Do NOT upload refs.bib to
%% Overleaf: it would only add this extra pass.
\IfFileExists{refs.bib}{\PassOptionsToPackage{pagebackref}{hyperref}}{}
\usepackage[breaklinks,colorlinks,allcolors=cvprblue]{hyperref}

\def\confName{CVPR}
\def\confYear{2027}

\title{Memory Is Not Enough: Separating Drift from Background Conflict\\in Continual Document Understanding}

\author{Anonymous CVPR submission}

%% Headline numbers: generated from the result files (never typed by hand).
\input{vals}
%% Merged view: the appendix (supp_body.tex) follows the references in this PDF.
%% CVPR 2027 kit practice: compile together, then split. make_submission.sh writes
%% submission/main_paper.pdf (pages up to the references) and submission/supplementary.pdf.
%% With \withappendixfalse the body's \cref{app:...}, \cref{fig:heatmaps} and \cref{eq:infer}-style
%% pointers into the appendix print as ?? (accepted: build supp.tex alongside; the merged view is the default).
\newif\ifwithappendix \withappendixtrue
\newcommand{\bodyref}[2]{\cref{#1}}  % appendix -> body pointer (supp.tex prints #2)
\ifwithappendix\input{supp_vals}\input{supp_perclass_vals}\input{image_vals}\fi

\begin{document}
\maketitle

\begin{abstract}
We study catastrophic forgetting issues in document understanding as key-information extraction (KIE) models must adapt to new documents arriving in a streaming fashion, which often leads them to forget what they learned from previous ones. Whereas replaying stored documents shows competitive performance against joint training when new domains arrive, it suffers from a big gap when the label space grows, known as class-incremental learning (CIL). We investigate the cause of this gap. Replay can reduce drift while retaining conflicting supervision through the background label, which marks every token that is not an annotated entity: a shared softmax treats each background label as evidence against all classes, while it is evidence only against the classes its annotator could see. Based on this finding, we propose a novel method called CoLaR-SBR. Inspired by latent replay and per-session readouts, our method maintains a Compressed Latent Replay (CoLaR) memory of whole documents to remind the model of past sessions, and uses a Session-Blocked Readout (SBR) that gives each session its own background logit to remove this conflict. We also show that per-session low-rank updates can replace the memory when all sessions share one annotation schema. We benchmark our method against various continual learning baselines on 5 public document datasets and achieve significant improvement in the class-incremental setting, with competitive performance against joint training.
\end{abstract}

\section{Introduction}
\label{sec:intro}

The recent literature has shown the high potential of multimodal encoders such as LayoutLMv3~\citep{huang2022layoutlmv3}, LiLT~\citep{wang2022lilt} and BROS~\citep{hong2022bros} in key-information extraction (KIE), where each OCR token of a visually-rich document (form, invoice, receipt) is labelled using its words, its 2-D layout and the page image. However, when documents arrive in a streaming fashion, with new document types and growing label sets, these models suffer from catastrophic forgetting~\citep{mccloskey1989catastrophic}. Retraining on all history is often blocked by cost and by privacy constraints. Thus, a deployed KIE model has to learn continually.

There exist many continual learning (CL) methods to cope with catastrophic forgetting, and the strongest ones in the class-incremental setting keep stored samples~\citep{masana2022class,wang2024survey}. On documents, replay looks sufficient at first. When new domains arrive under one fixed label space, i.e., domain-incremental learning (DIL), replaying a few hundred stored documents keeps LayoutLMv3 close to joint training (\cref{app:drift}). Unfortunately, when the label space grows session by session, i.e., class-incremental learning (CIL), replay ends far below joint training (\cref{tab:main}). A class-incremental stream is expected to be harder. Nevertheless, it is unclear why a memory that holds the old documents together with their labels does not help. In KIE, every token that is not an annotated entity receives the tag ``O'', which we call the background label. The shift of this label is known from class-incremental segmentation~\citep{cermelli2020mib,douillard2021plop}, incremental detection~\citep{liu2023abr,mo2024bpf} and continual named-entity recognition (NER)~\citep{monaikul2021extendner,zhang2023cpfd}, and separated softmaxes~\citep{ahn2021ssil,kim2022more,liu2023ice} already undo the shared normalisation. However, to our knowledge, no prior work examines, on token-labelled documents, which part of the failure a memory repairs and which part remains.

\input{fig/conflict}

In this work, we investigate the cause of this gap and propose a novel method called CoLaR-SBR. First, we examine a series of hypotheses on class-incremental streams (\cref{tab:ladder}). For each hypothesis, we state what it predicts, run an intervention on the same encoder, stream and recipe, and measure where the errors go. We find that label growth does not only make forgetting harder: with a whole-document memory, the model suppresses the session it is learning, which naive training does not do. Then, on CIL-CORD, an oracle that relabels every stored document with its true labels brings the shared softmax close to joint training, whereas on CIL-FUNSD and CIL-WildReceipt it does not, and on CIL-FUNSD a loss that lets the background label mean ``background or an earlier class''~\citep{cermelli2020mib} does. Giving each session its own background closes most of the gap on all three class-incremental streams (\cref{fig:repairs}). These results suggest that replay can reduce drift while retaining conflict in supervision through the background label, since a shared softmax treats every background label as evidence against all classes, while it is evidence only against the classes its annotator could see (\cref{fig:conflict}).

Based on this finding, we propose CoLaR-SBR (\cref{fig:method}). Our method maintains a Compressed Latent Replay (CoLaR) memory, inspired by latent replay~\citep{hayes2020remind,pellegrini2020latent}, which stores the compressed activations of whole documents together with their labels and the session that wrote them, to remind the model of past sessions. On top of that, we use a Session-Blocked Readout (SBR), inspired by ICE-O~\citep{liu2023ice} and MORE~\citep{kim2022more}, which gives each session its own background logit and its own softmax. We argue that each class should be supervised only with negatives from the documents that annotated it, as in training on disjoint tag sets~\citep{greenberg2018marginal}. Hence, a background label trains only the blocks of the sessions that annotated the document. Furthermore, we show that the memory can be replaced in one case: when all sessions share one annotation schema, per-session low-rank updates~\citep{hu2022lora} (SBR-Recall) come close to joint training with no stored document.

We benchmark our method against two references and eight CL baselines on six streams over five public document datasets, FUNSD, CORD, SROIE, WildReceipt and XFUND. On the three class-incremental streams, CoLaR-SBR closes most of the gap to joint training on every seed. It is noted that our readout is not limited to the latent memory: with raw documents instead, the full-label F1 differs by at most $\nErSbrAbsFullMax$ points on every seed. The contributions of our study are as follows: (1) we investigate why replay fails under label growth with controlled interventions and find that replay can reduce drift while retaining background-label conflict, with different repairs succeeding on different streams (\cref{sec:investigation}); (2) we propose a novel method called CoLaR-SBR, which combines a whole-document latent memory with a per-session background readout to address this conflict without additional annotations, and we show when per-session low-rank updates can replace the memory (\cref{sec:design}); (3) we conduct extensive experiments against eight CL baselines on six streams over five public document datasets, where our method achieves significant improvement on the three class-incremental streams (see \cref{tab:main}); our ablation study further shows that the memory and the readout address two different causes and are both needed, that the gain comes from the readout and not from the format of the memory, and that it holds with other class orders and with a second backbone (\cref{sec:beyond,sec:ablations}).
\section{Related Work}
\label{sec:related}
\paragraph{Document understanding and continual document IE.} Recent studies on visually-rich documents mainly consist of layout-aware encoders~\citep{xu2020layoutlm,xu2021layoutlmv2,huang2022layoutlmv3,wang2022lilt,hong2022bros} and generative models~\citep{kim2022donut,tang2023udop}. They achieve strong results on forms and receipts. However, they are trained once on all the data, with a classifier sized to a fixed label set. There have been few works on continual learning for documents: continual table detection uses plain ER~\citep{minouei2022table}, and ProtoNER adds new classes with prototypes~\citep{kumar2023protoner}. Recently, \citet{sassioui2026documents} benchmarked the forgetting of layout-aware models under domain, layout and schema shifts. However, they decompose forgetting by the kind of shift, whereas we decompose it by cause. It is noted that a generative model emits a text sequence and no per-token label, so it has no explicit background class; hence, we work with token classifiers (\cref{sec:limitations}).

\paragraph{Continual learning and replay.} Recent studies on catastrophic forgetting mainly consist of regularization~\citep{kirkpatrick2017ewc,li2017lwf}, parameter isolation~\citep{mallya2018packnet,rusu2016pnn} and replay~\citep{rolnick2019er,buzzega2020der}, and the best methods in a large class-incremental comparison keep exemplars~\citep{masana2022class}. More recently, methods built on a frozen pre-trained model learn prompts~\citep{wang2022l2p,smith2023coda} or low-rank adapters~\citep{liang2024inflora,he2025cllora}, or rebuild the classifier from class statistics~\citep{zhang2023slca,mcdonnell2023ranpac}. These methods are mainly designed for image classification, where each sample carries one label and no class plays the role of a background. To reduce the cost of raw samples, latent replay stores intermediate activations and freezes~\citep{hayes2020remind} or slows down~\citep{pellegrini2020latent} the layers below them. Whereas most prior replay mechanisms focus on which samples to store or how to compress them, we investigate another aspect, namely what the stored labels say: under label growth, a stored document carries the background label on every entity outside its session.

\paragraph{Background shift.} When each step annotates only its own classes, the background label changes meaning. In class-incremental segmentation, MiB~\citep{cermelli2020mib} lets a background label mean background or any old class in the loss, PLOP~\citep{douillard2021plop} pseudo-labels the background with the old model, and SSUL~\citep{cha2021ssul} models an unknown class inside the background, with per-class sigmoid outputs on a frozen backbone. \citet{kalb2022causes} find that the semantic shift of the background and the bias towards new classes are the main causes of forgetting there. Incremental detection has named both the stored side of the shift~\citep{liu2023abr} and its past-and-future form~\citep{mo2024bpf}. In continual NER, it is the unlabelled-entity problem~\citep{ma2023learning}, answered by distillation from the old model on background tokens~\citep{monaikul2021extendner,zheng2022cfner}, pseudo-labels~\citep{zhang2023cpfd}, and distillation with a debiased loss~\citep{qiu2024is3}. These methods name the shift, and many of them answer it through the old model. We do not claim the conflict; instead, we separate what a replay memory repairs from what survives true stored labels, and measure what each repair leaves.

\paragraph{Separated readouts and task identity.} A softmax over old and new classes together biases their scores, and separating it, with task-wise distillation, reduces the bias~\citep{ahn2021ssil}; ER-ACE~\citep{caccia2022erace} restricts the softmax of incoming samples to their own classes. Class-incremental learning decomposes into within-task prediction and task identification, which is an out-of-distribution problem~\citep{kim2022theoretical}; a ``none of mine'' output per task is one way to solve it. MORE~\citep{kim2022more} isolates each task's parameters and trains an out-of-distribution output per head from saved samples, and ICE-O~\citep{liu2023ice}, on a frozen encoder with no memory, trains one softmax per session over a constant Other logit on text information extraction; its replay variant stores samples per class and tunes the encoder. Unlike them, SBR learns one background per session and keeps old blocks trainable from a memory, and SBR-Recall shares MORE's isolation, by low-rank deltas, with no saved samples. To our knowledge, on token-labelled documents, no prior work combines a memory against drift with one learned background per session, and none measures what each one repairs.

\section{Problem Formulation}
\label{sec:setup}
In the following, we first define continual key-information extraction and the three settings we study. Then, we describe how the streams of each setting are built, and how a model is scored.

\paragraph{Continual key-information extraction.} A document is its OCR tokens, their normalised boxes and the page image. The model labels every token with a BIO tag, and the entities are scored by entity-level F1 (\texttt{seqeval}). The documents arrive in $T$ sessions. Session $t$ brings a training set $\mathcal{D}_t$ that is annotated with its own tag set $\mathcal{T}_t$ only: a token whose true tag is $y^{\ast}$ is labelled
\begin{equation}
y=\begin{cases} y^{\ast} & y^{\ast}\in\mathcal{T}_t\\ \mathrm{O} & \text{otherwise.}\end{cases}
\label{eq:mask}
\end{equation}
In the BIO scheme, ``O'' is the tag of every token that belongs to no annotated entity. Hereinafter, we call it the \emph{background label}: we write ``O'' when we refer to the tag in the data, and background label when we refer to its role in training. Under \cref{eq:mask}, the background label of session $t$ covers both the tokens that are truly not entities and the entities that session $t$ does not annotate. The model is trained on $\mathcal{D}_1,\dots,\mathcal{D}_T$ in turn and cannot return to an earlier session, except through what its method stores. After the last session, it is tested on the classes of all sessions, without knowing which session a document comes from.

\paragraph{Three settings.} We study three settings, six streams in total, over five public document datasets: FUNSD~\citep{jaume2019funsd}, CORD~\citep{park2019cord}, SROIE~\citep{huang2019sroie}, WildReceipt~\citep{sun2021sdmgr} and XFUND~\citep{xu2022xfund} (\cref{tab:bench}). Each setting has its own role. In \emph{class-incremental learning} (CIL), the label set grows and the documents stay in one domain; this is where the background label changes meaning, hence where our main claims are tested. In \emph{domain-incremental learning} (DIL), the label set is fixed and the documents change; no label is masked there, so forgetting can only come from the drift of the model, and DIL serves as the control without the conflict. In the \emph{Mixed} setting both change, which tests the two causes together.

\paragraph{Class-incremental streams.} We take CORD~\citep{park2019cord} as the running example (\cref{tab:cordsessions}). CORD has $800$ training receipts and $100$ test receipts, each annotated with $30$ entity classes such as the name and the price of a menu item, the subtotal, the tax and the total. We do not change the receipts; we change which of their labels a session can see. First, the $30$ classes are divided into five groups of six, one group per session. Then, the training set of session $t$ keeps every receipt that holds at least one entity of its group, and applies \cref{eq:mask}: the six classes of the session keep their labels, and every other entity on the receipt, of an earlier or of a later session, is labelled ``O''. Finally, $10\%$ of this training set is held out for early stopping. Three consequences follow. The classifier head grows by the tags of each new session. The same receipt returns in several sessions under different label masks, as in the overlapped protocol of class-incremental segmentation~\citep{cermelli2020mib}. And the same token receives contradicting labels over time: the total of a receipt is ``O'' in session~1 and \texttt{total.total\_price} in session~5, whereas the price of a menu item is \texttt{menu.price} in session~1 and ``O'' in session~5 (\cref{fig:conflict}). Joint training, our reference, sees every receipt once with all of its $30$ labels.

\begin{table}[t]
\centering\footnotesize
\setlength{\tabcolsep}{3pt}
\caption{\textbf{The five sessions of CIL-CORD.} Each session annotates six of CORD's $30$ classes; every other entity is labelled ``O''. Receipts: the CORD training receipts that hold at least one entity of the session ($800$ in all).}
\label{tab:cordsessions}
\resizebox{\linewidth}{!}{\begin{tabular}{@{}clc@{}}
\toprule
Session & Classes annotated (others are ``O'') & Receipts \\
\midrule
1 & menu: cnt, discountprice, itemsubtotal, nm, num, price & 800 \\
2 & menu: unitprice, vatyn, sub.cnt, sub.nm, sub.price, sub.unitprice & 382 \\
3 & sub\_total: discount, etc, othersvc, service, subtotal, tax & 549 \\
4 & total: cashprice, changeprice, creditcardprice, emoneyprice, menuqty\_cnt, menutype\_cnt & 723 \\
5 & total: total\_etc, total\_price; void\_menu: nm, price; sub: nm, cnt & 780 \\
\bottomrule
\end{tabular}}
\end{table}

The three CIL streams are built in this way. CIL-CORD has five sessions of six classes. CIL-FUNSD has two sessions on the $149$ training forms of FUNSD: the first annotates headers and questions, the second annotates answers. CIL-WildReceipt has four sessions of six classes on the $1{,}267$ training receipts of WildReceipt, with its $24$ classes taken in the order of the dataset.

\paragraph{Domain-incremental streams.} Here every dataset is mapped to one schema of four classes (header, key, value and other) before training. For instance, a question of a FUNSD form becomes a key and its answer a value, and the price or the date of a receipt becomes a value. Every session is annotated with all four classes, so \cref{eq:mask} masks nothing and the head never grows. Each session is a different set of documents, and the only shift is the domain or the language. In DIL, the model learns FUNSD forms, then SROIE receipts, then CORD receipts. In DIL-XFUND, it learns the XFUND forms in German, Spanish, French, Italian and Chinese in turn ($149$ training forms per language).

\paragraph{Mixed stream.} Mixed has six sessions over three datasets and their three annotation schemas: (1) FUNSD headers and questions, (2) FUNSD answers, (3) the four fields of SROIE (company, date, address and total), (4) and (5) two groups of CORD classes, and (6) FUNSD again with all of its labels and no new class. Hence, a boundary brings new classes (1 to 2, 4 to 5), a new domain with a new schema (2 to 3, 3 to 4) or a return to an old domain (5 to 6). The head grows as in CIL, and the entities of another session are labelled ``O'' (\cref{app:setup}).

\begin{table*}[t]
\centering\scriptsize
\setlength{\tabcolsep}{2.5pt}
\caption{\textbf{Benchmark settings}: six streams over five document datasets (sessions: \cref{app:setup}). Train/test documents: FUNSD 149/50, SROIE 626/347, CORD 800/100. AA: mean F1 of the final model over the sessions; full-label F1: on the unmasked test set of the stream's dataset (\cref{sec:setup}).}
\label{tab:bench}
\input{tab_bench}
\end{table*}

\paragraph{Evaluation protocol.} After the last session, the final model is scored in two ways. Our headline number on the class-incremental streams is the \emph{full-label F1}: the entity-level F1 on the complete test set of the dataset with all of its classes and no label masked ($100$ receipts and $30$ classes on CIL-CORD), as in continual NER~\citep{ma2023learning}. To measure retention, we also keep one test split per session, built like its training set and masked in the same way. From $R[i,j]$, the F1 on the test split of session $j$ after training session $i$, we report the average accuracy AA $=\tfrac1T\sum_j R[T,j]$ and the average forgetting AF $=\tfrac1{T-1}\sum_{j<T}\big(R[j,j]-R[T,j]\big)$, the F1 of a session right after it was learned minus its final F1. Backward transfer (BWT) is the negative of AF under this definition, so we do not report it separately. We do not report forward transfer either: with a growing head, a class has no output before its session, so the model cannot be scored on a session it has not yet learned. Since the labels of a session split show only the classes of that session, raw scoring would charge a model for every correct entity of another session. Hence, we score each session split \emph{session-restricted}: a predicted tag outside the label set of session $j$ is mapped to ``O'' before span matching (\cref{app:setup}). A fixed label space has no masked split and no growing head, so DIL and DIL-XFUND report AA and AF only, and so does Mixed.

\section{Why Replay Fails under Label Growth}
\label{sec:investigation}
Every hypothesis is tested on CIL-CORD with LayoutLMv3 under the protocol of \cref{sec:protocol} and changes one intervention against the run it is compared with, unless stated. \Cref{tab:ladder} lists the hypotheses and \cref{fig:forgetting} shows the forgetting session by session. For each hypothesis, we report whether its prediction holds on the three seeds: it holds on every seed (\emph{supports}), fails on every seed (\emph{refutes}) or holds in part (\emph{partly}). It is noted that this is a comparison of signs over three seeds, not a significance test. The \emph{blocked readout} (\cref{sec:sbr}) gives each session's tags their own background logit and softmax, a \emph{block}. Unless stated, a memory is CoLaR (\cref{sec:memory}): $50$ documents per session stored as rank-$128$ SVDs of their layer-$4$ activations.

\begin{table*}[t]
\centering\footnotesize
\setlength{\tabcolsep}{3pt}
\caption{\textbf{The series of hypotheses} on CIL-CORD with LayoutLMv3, unless stated (AF: average forgetting, mean over seeds; other cells: mean $\pm$ s.d. over three seeds; $^{n=k}$: $k$ seeds; $^\dagger$legacy protocol). $^{r}$Oracle: stored documents relabelled with their true labels (\cref{sec:h1}). $^{e}$ER trains every layer and keeps $\nErCilBuf$ raw documents. The pre-session row of H1 also tests H2.}
\label{tab:ladder}
\input{tab_ladder}
\end{table*}

\begin{figure}[t]
\centering
\includegraphics[width=\linewidth]{figures/repairs.pdf}
\caption{\textbf{Three repairs on three class-incremental streams} (full-label F1; bar: mean of three seeds, dots: seeds; dashed: joint training). Every repair trains with the same CoLaR memory, and the first three keep the shared softmax. True stored labels, an oracle, close the gap on CIL-CORD only and the MiB-style loss on CIL-FUNSD only, whereas the per-session readout closes most of it on all three streams (values: \cref{tab:repairs}).}
\label{fig:repairs}
\end{figure}
\begin{figure*}[t]
\centering
\includegraphics[width=\textwidth]{figures/forgetting.pdf}
\caption{\textbf{Forgetting on CIL-CORD}, session by session (line: mean of three seeds, band: their range). (a) F1 averaged over the sessions seen so far. (b) F1 on session~1. Dashed purple: the memory with true labels before each session, an oracle. Dotted: joint training, one pooled run.}
\label{fig:forgetting}
\end{figure*}

\subsection{H0: only harder forgetting}
\label{sec:h0}
We first hypothesize that label growth only makes drift harder. Then the memory that retains under domain shift would still retain here, and the losses would look like naive forgetting, where old sessions are lost and the newest is learned well. However, the results refute it. The CoLaR memory comes within $\nColarDilGapMax$ AA points of joint training on DIL on every seed, whereas on CIL-CORD with a shared softmax it reaches only $\vColarFull$ full-label F1 against $\vJointFull$, and raw documents do no better (ER with $\nErCilBuf$ documents: $\vErFull$). What separates the memory run from naive forgetting is the session being trained: its newest session ends at $\mColarLastMin$ to $\mColarLastMax$ F1, whereas naive training learns every session to at least $\nNaiveDiagMin$ on every seed (\cref{app:results}). With the same frozen layers and the same replay, the blocked readout learns every session to at least $\nSbrDiagMin$, so lost plasticity does not explain it. Hence, drift alone does not explain the failure. Nevertheless, drift is real: on a fixed label space every encoder forgets its earlier domains, and replay repairs it (\cref{app:drift}).

\subsection{H1: stale labels in the memory}
\label{sec:h1}
We next hypothesize that the stored labels are stale. A document stored at session $t$ carries the background label on every entity of another session, so the memory itself calls background what the model must keep or learn; in detection, this staleness was named foreground shift~\citep{liu2023abr}. If so, true stored labels should close the gap with the shared softmax unchanged. To examine it, we repair the stored labels with two oracles, which read labels no continual learner has. First, \emph{before each session}, every stored document receives the true labels of all classes in the grown head, so no stored label is stale during training, while the current documents still label old entities as background. Second, we mask the same stale tokens in replay instead, removing the false background label without adding a label (a refresh at the session boundaries only: \cref{app:results}).

The hypothesis is supported on CIL-CORD. Refreshed before each session, the same memory and softmax reach $\vOraclePreFull$ full-label F1, within $\nOraclePreJointGap$ points of joint training on every seed. Masking instead reaches $\vOracleMaskFull$; hence, removing the stale negatives alone recovers most of the gap. Thus, repairing stored supervision recovers nearly all of the gap on CIL-CORD; relabelling also adds positive labels, while masking isolates removal of false negatives. However, this does not hold on every stream: on CIL-FUNSD the same oracle leaves the shared softmax where it was, and on CIL-WildReceipt it helps substantially on one seed only (\cref{fig:repairs}).

\subsection{H2: the background labels of the current documents}
\label{sec:h2}
We then hypothesize that the current documents are the cause: each labels every entity of an earlier session as background, so even a memory with true labels is opposed in every batch; this is the second side of the conflict named in detection~\citep{mo2024bpf}. If so, a loss that lets a current background label mean ``background or an earlier class'', as MiB does for segmentation~\citep{cermelli2020mib} and marginal-likelihood taggers do offline~\citep{greenberg2018marginal}, should close the gap within a shared softmax. To examine it, we train the memory and the shared softmax with two repairs requiring no additional annotations. The first is MiB's unbiased cross-entropy in both directions: a current background token matches the background or any earlier-session class, and a replayed one the background or any class introduced after storage (no distillation). The second slices each stored document's replay loss to the classes known at storage.

The results show that a loss repair only partly helps on CIL-CORD. The MiB-style loss lifts full-label F1 to $\vMibFull$, \nMibShareFull\% of the seed-matched gap, and the softmax slice reaches only $\vSlicedFull$ (\cref{app:results}). Besides, the oracle of H1 closes the gap although every current document still labels old entities as background; thus, current-side conflict does not prevent rescue by this oracle on CIL-CORD. On CIL-FUNSD, the bidirectional MiB-style loss is level with our readout; this intervention alone does not isolate the current side. On CIL-WildReceipt, it recovers more than the oracle but stays below our readout on every seed. Besides, slicing replay and current-background losses, a partial-label control inspired by SS-IL~\citep{ahn2021ssil} without distillation, reaches only $\vSlicedBothFull$ on CIL-CORD, and slicing the current-background loss alone $\vSlicedCurFull$, both at least $\nSlicedAnyBelowSbrMin$ points below our readout on every seed.

\subsection{H3: the shared normalisation}
\label{sec:h3}
We hypothesize that the conflict lives in the shared normalisation, one background logit trained by documents that disagree on the background. Then one background logit per session, normalised with that session's tags only, should close most of the gap without additional annotations. This is SS-IL's diagnosis of a shared softmax~\citep{ahn2021ssil} with a background added, and ICE-O's readout~\citep{liu2023ice} without its frozen encoder. We test it with and without memory to measure the contribution of replay.

The hypothesis is supported on every seed. With the memory, the blocked readout (CoLaR-SBR, \cref{sec:sbr}) reaches $\vSbrFull$ full-label F1, $\nSbrGapFullMin$ to $\nSbrGapFullMax$ points below joint training per seed. Without the memory, the same readout scores $\vNoBankFull$: old blocks on a drifted trunk claim far more than their own entities (\cref{app:results}), consistent with drift and readout miscalibration; replay reduces this error. \Cref{fig:repairs} compares the three repairs on the three class-incremental streams: the label-side repairs each succeed on one stream, whereas the blocked readout closes most of the gap on all three. The readout is not free: on DIL, with nothing to separate, CoLaR-SBR is slightly below CoLaR (\cref{tab:main}).

\section{Method}
\label{sec:design}
In this section, we describe CoLaR-SBR (\cref{fig:method}). In the following, we first state the principle behind it. Then, we detail its memory (\cref{sec:memory}), its readout (\cref{sec:sbr}) and a variant that replaces the memory (\cref{sec:recall}).

\paragraph{The principle.} These results motivate a supervision rule under label growth, familiar from training on partially annotated data~\citep{greenberg2018marginal}. A token labelled as background in session $s$ means ``none of the classes annotated in session $s$'': it is evidence against those classes and silent about every other. However, a shared softmax turns it into evidence against all classes, and a stale stored label makes the same error inside the memory. We therefore test the design principle: \emph{supervise each class only with negatives from documents that annotated it, and normalise together only classes that were annotated together.} We instantiate this principle with one background per session and whole-document memory retaining labels and annotation-session identity; and each session's background logit, trained on that session's true negatives, is also a detector of its entities, so it can route a token without task identity. We name the resulting method CoLaR-SBR. \Cref{tab:scope} reads every repair of \cref{sec:investigation} through this principle: they change the scope of negative supervision, but also positive labels, loss geometry or routing. Their success supports the principle as a repair, not as a necessary condition.

\begin{table}[t]
\centering\footnotesize
\setlength{\tabcolsep}{3pt}
\caption{\textbf{One principle, six repairs} (CIL-CORD, same CoLaR memory, full-label F1). How each repair treats a stored background label. $^{r}$Oracle: reads the true labels.}
\label{tab:scope}
\resizebox{\linewidth}{!}{\begin{tabular}{@{}llc@{}}
\toprule
Repair & Treatment of stored background label & F1 \\
\midrule
Shared softmax & background against all entities & $\vColarFull$ \\
Softmax slice & background against capture-time entities & $\vSlicedFull$ \\
MiB-style loss & marginal target: ``O or a later class'' & $\vMibFull$ \\
Stale tokens masked$^{r}$ & no loss where the label is stale & $\vOracleMaskFull$ \\
True stored labels$^{r}$ & corrected to the true label & $\vOraclePreFull$ \\
Per-session readout (SBR) & background within its own session & $\vSbrFull$ \\
\bottomrule
\end{tabular}}
\end{table}

\input{fig/method}

\subsection{CoLaR: a whole-document latent memory}
\label{sec:memory}
Our readout (\cref{sec:sbr}) needs stored documents that keep each token's label at its position and the session that wrote it. CoLaR (Compressed Latent Replay), inspired by latent replay~\citep{hayes2020remind,pellegrini2020latent}, keeps whole documents: after the first session layers $1..k$ are frozen, and each stored document's layer-$k$ activations $H_j$ are kept as a rank-$r$ truncated SVD with its boxes, mask, labels and session, rebuilt at replay and injected at layer $k$; it holds no raw pixels or token ids (no privacy claim: activations can leak content). The stored documents are drawn at random from the session's training set. We use CoLaR because it keeps no pixels or token ids (\cref{sec:cost}).

\subsection{SBR: one background per session}
\label{sec:sbr}
SBR (Session-Blocked Readout) is our instantiation of a per-session background readout~\citep{liu2023ice,kim2022more} (\cref{fig:method}b). Every session that brings new tags opens a block $b$ holding those tags $\mathcal{T}_b$. Whereas a shared softmax normalises over $\{\mathrm{O}\}\cup\bigcup_b\mathcal{T}_b$, SBR gives every session its own background logit $o_b(x)=w_b^{\top}h(x)+c_b$ on the token feature $h(x)$ and normalises each block separately:
\begin{equation}
p_b(\cdot\mid x)=\operatorname{softmax}\big([\,o_b(x),\,z_{\mathcal{T}_b}(x)\,]\big),
\label{eq:block}
\end{equation}
where $z(x)$ are the tag logits of the usual linear head and $z_{\mathcal{T}_b}(x)$ its entries for the tags in $\mathcal{T}_b$. Within block $b$ a token labelled ``O'' in session-$b$ data is truly background, so the block's supervision is never stale. For a token $x$ with label $y$ in a document of session $s$,
\begin{equation}
\begin{gathered}
\mathcal{L}(x,y,s)=\sum_{b}\,m_{sb}\;\mathrm{CE}\big(p_b(\cdot\mid x),\,y_b\big),\\
y_b=\begin{cases} y & y\in\mathcal{T}_b\\ \mathrm{O}_b & \text{otherwise,}\end{cases}
\end{gathered}
\label{eq:sbr}
\end{equation}
with one rule for the mask: $m_{sb}{=}1$ when session $s$ annotates every tag of $\mathcal{T}_b$; $m_{sb}{=}1$ with $y_b{=}\mathrm{O}_b$ when $b$ belongs to a different annotation schema (every token is a negative for $b$, a training-time convention; \cref{app:method}); and $m_{sb}{=}0$ otherwise, because the document's background labels say nothing about block $b$. Stored documents follow the rule of the session that stored them. A shared softmax is the special case of a single block that every session supervises ($m_{sb}{=}1$ for all $s$), which is where a stale background label becomes a wrong negative. Since a session that adds no tags creates no block, SBR has a single block on a fixed label space. At inference each block reports its claim $e_b=1-p_b(\mathrm{O}_b\mid x)$; the block with the largest $e_b$ labels the token with its best tag if $e_b>0.5$, and the token is ``O'' otherwise, as in ICE-O's ``Other unless some session scores above it''. Blocks trained at different times need not produce claims on one scale; hence we measure the routing directly (\cref{fig:findings}b).

\subsection{SBR-Recall: per-session deltas instead of the memory}
\label{sec:recall}
We hypothesize that, without the conflict, the memory is needed only against drift. The task-identity view~\citep{kim2022theoretical} says when it can go: if each block is read under the trunk state that learned it, nothing drifts for it, provided the model can tell which state a token needs. On one schema under label growth each block's background logit supplies that key. Across schemas, however, the all-background terms of \cref{eq:sbr} train every old block towards background on other-schema documents with nothing to pull it back, so isolation should fail without stored documents; a fixed label space has one block and nothing to route by, so the deltas should fail there as well. To examine this hypothesis (H4), SBR-Recall freezes the layers above $k$ after the first session (the base $\theta_0$), stores for each later session $b$ a rank-$64$ delta $\Delta_b$~\citep{hu2022lora} of them ($\Delta_1{=}0$), and computes every term of \cref{eq:sbr} for block $b$, at training and test time, under the trunk state that learned it (\cref{app:method}):
\begin{equation}
p_b(\cdot\mid x)=\operatorname{softmax}\big([\,o_b,\,z_{\mathcal{T}_b}\,]\big)\big(x;\,\theta_0+\Delta_b\big).
\label{eq:recall}
\end{equation}
We report its results in \cref{sec:ablations}.

\section{Experiments}
\label{sec:experiments}
\subsection{Experimental setup}
\paragraph{Model and input.} Every run fine-tunes LayoutLMv3-base~\citep{huang2022layoutlmv3} with one linear token-classification head, which grows by the tags of each new session on the CIL and Mixed streams and is fixed on the DIL streams. We use the OCR words and boxes provided with each dataset, page images at $224{\times}224$ and at most $512$ text tokens per document.

\paragraph{References.} We report two references. \emph{Naive fine-tuning} trains on the sessions in turn with no continual-learning mechanism, which is the lower reference. \emph{Joint training} trains once on all documents of the stream with all of their labels, which is the reference we compare the gap against.

\paragraph{Baselines.} We benchmark against eight methods from five families (\cref{tab:main}). \emph{Replay:} ER~\citep{rolnick2019er} keeps a reservoir of raw documents with their session labels and trains on them together with the current session, and DER++~\citep{buzzega2020der} also replays the logits stored with them. \emph{Regularization:} LwF~\citep{li2017lwf} distils the outputs of the previous model and stores no document. \emph{Prompts and low-rank adapters on a frozen encoder:} CODA-Prompt~\citep{smith2023coda} and CL-LoRA~\citep{he2025cllora}. \emph{Continual NER:} CPFD~\citep{zhang2023cpfd} and IS3~\citep{qiu2024is3}, which are designed for the shift of the background label. \emph{Merging:} LCA~\citep{lca2026}, which merges per-session backbones and realigns the classifier. All of them are our ports onto the one growing token head that every row uses, and all of them keep the shared softmax (\cref{app:baselines}).

\paragraph{Our methods.} \emph{CoLaR} is the memory alone, with a shared softmax. \emph{CoLaR-SBR} adds the session-blocked readout, and \emph{SBR-Recall} replaces the memory with per-session low-rank deltas (\cref{sec:design}).

\paragraph{Memory.} A memory is filled at the end of each session with documents of that session, stored with the masked labels they had in that session. CoLaR and CoLaR-SBR keep $50$ documents per session ($10$ on CIL-FUNSD, which has $149$ training forms), i.e., $250$ documents after the five sessions of CIL-CORD, as rank-$128$ SVDs of their layer-$4$ activations; the embeddings and layers $1$ to $4$ are frozen after the first session. ER and DER++ keep raw documents, with the same budget as ours on CIL-CORD ($250$) and on CIL-FUNSD ($20$). All memory methods replay four stored documents per current batch of two.

\paragraph{Training protocol.}\label{sec:protocol} We use AdamW with learning rate $5{\times}10^{-5}$ (LCA: its own SGD recipe), full precision and batch size $2$, which is a hardware limit. Every session is trained with early stopping on the $10\%$ of its \emph{training} set that was held out: training stops when the F1 on this held-out part has not improved for two epochs, and the weights of the best epoch are restored. Hence, the test split never selects a model. Joint training follows the same rule, with a patience of ten epochs on CIL-CORD and CIL-FUNSD (\cref{app:setup}). Every cell is the mean and the standard deviation over three seeds (42, 7 and 123), which change the initialisation of the head, the order of the batches and the stored documents, but not the sessions. Cells marked $\dagger$ come from earlier runs that stopped early on the test split (\emph{legacy protocol}); they are given for completeness (\cref{app:setup}).

\paragraph{Hyperparameter tuning.} Tuning and evaluating on the same scenario overstates CL methods~\citep{cha2025hyperparameters}. Hence, the comparison set was fixed after a first grid on all three streams, the memory settings were chosen on DIL, the claim threshold of SBR ($0.5$) was not tuned, and the one setting chosen on CIL-CORD is the delta capacity of SBR-Recall (\cref{app:baselines,app:results}). The baselines use the default hyperparameters of their methods.

\begin{table*}[t]
\centering\small
\setlength{\tabcolsep}{4pt}
\caption{\textbf{Comparison with continual-learning baselines} on LayoutLMv3 over six streams. F1: full-label F1; AA and AF: average accuracy and forgetting (\cref{sec:setup}). AA and F1: mean $\pm$ s.d. over three seeds, per-seed values when fewer; AF: mean; ``--'': not run. Bold: best of each column, Joint training aside; underlined: our best when a baseline is best. $^\dagger$legacy protocol; $^{n=k}$: $k$ clean seeds so far; $^{b}$with the $50$-document memory; on CIL-WildReceipt and DIL-XFUND, ER keeps $200$ documents and DER++ $250$; on CIL-FUNSD both keep $20$, the memory of the CoLaR cells there; Ports and protocols: \cref{app:baselines}; what each method stores: \cref{tab:memcmp}.}
\label{tab:main}
\resizebox{\linewidth}{!}{\input{tab_main_clean}}
\end{table*}

\subsection{Results}
\label{sec:beyond}
\paragraph{Results on CIL-CORD.} \Cref{tab:main} shows the comparison between our method and the baselines. CoLaR-SBR reaches $\vSbrFull$ full-label F1 ($\vSbr$ AA), against $\vJointFull$ for joint training and $\vColarFull$ for the same memory with a shared softmax, and SBR-Recall reaches $\vRecallFull$ with no stored document. Every shared-softmax method of \cref{tab:main} stays far below them: the best baseline, ER with $\nErCilBuf$ raw documents, reaches $\vErFull$ full-label F1, below CoLaR with a shared softmax. It is noted that several of our ports learn the later sessions poorly, and more patience does not close the gap: with a patience of ten epochs instead of two on CIL-FUNSD, the four weakest ports change by \nProbeDeltaMin{} to \nProbeDeltaMax{} AA points, and the best of them stays $\nProbeBestGapSbr$ points below CoLaR-SBR (seed 42; \cref{app:baselines}). It does not depend on the class order either: with the sessions reversed and shuffled (seed 42), CoLaR reaches $\nOrderRevColar$ and $\nOrderShufColar$, CoLaR-SBR $\nOrderRevSbr$ and $\nOrderShufSbr$, and SBR-Recall $\nOrderRevRecall$ and $\nOrderShufRecall$. Besides, the finding holds on a second backbone: with BROS, CoLaR-SBR reaches $\vBrosSbrFull$ full-label F1 against $\vBrosColarFull$ for CoLaR (\cref{app:results}).

\paragraph{Results on CIL-FUNSD and CIL-WildReceipt.} We observe the same pattern on the two other class-incremental streams. On CIL-FUNSD (two sessions, ten stored forms per session), CoLaR with a shared softmax stays close to naive fine-tuning ($\vFunsdColarFull$ against $\vFunsdNaiveFull$ full-label F1), whereas CoLaR-SBR reaches $\vFunsdSbrFull$ and SBR-Recall $\vFunsdRecallFull$, level with joint training ($\vFunsdJointFull$). Besides, every baseline of \cref{tab:main} stays at least $\nFunsdSbrMinusBestBaseMin$ AA points below CoLaR-SBR on every seed, LCA being the best of them ($\vFunsdBaseLca$ AA). On CIL-WildReceipt (four sessions), CoLaR with a shared softmax reaches $\vWildColarFull$ full-label F1, close to naive fine-tuning ($\vWildNaiveFull$), while CoLaR-SBR reaches $\vWildSbrFull$, which is \nWildSbrMinusColarFullMin{} to \nWildSbrMinusColarFullMax{} points above CoLaR on every seed and \nWildJointMinusSbrFullMin{} to \nWildJointMinusSbrFullMax{} points below joint training ($\vWildJointFull$). SBR-Recall, which stores no documents, reaches $\vWildRecallFull$.

\paragraph{Results on Mixed, DIL and DIL-XFUND.} On Mixed, where the sessions also change the annotation schema, CoLaR-SBR reaches $\vSbrMix$ AA, against $\vJointMix$ for joint training. On a fixed label space there is one block and nothing to separate. Hence, on DIL every replay method comes close to joint training and the blocked readout brings no gain (CoLaR-SBR $\nSbrDil$ against CoLaR $\nColarDil$ AA). On DIL-XFUND, CoLaR and CoLaR-SBR improve over naive fine-tuning, but they stay below ER and DER++, which keep raw documents and train every layer. As predicted in \cref{sec:recall}, SBR-Recall, which has no memory and nothing to route by on a fixed label space, falls below naive fine-tuning on both streams.

\paragraph{Forgetting analysis.} \Cref{fig:forgetting} shows how each method forgets on CIL-CORD. After five sessions, naive fine-tuning keeps $\nFirstNaive$ F1 on session~1 and CoLaR with a shared softmax $\nFirstColar$, whereas CoLaR-SBR keeps $\nFirstSbr$ and SBR-Recall $\nFirstRecall$. It is noted that CoLaR with a shared softmax has a low AF only because it suppresses each new session while learning it (\cref{sec:h0}); hence, we read AF together with AA.

\subsection{Ablation study}
\label{sec:ablations}
\begin{table}[t]
\centering\small
\setlength{\tabcolsep}{3pt}
\caption{\textbf{The two components} on CIL-CORD (mean $\pm$ s.d. over three seeds).}
\label{tab:2x2}
\resizebox{\linewidth}{!}{\input{tab_2x2}}
\end{table}

\paragraph{Ablation study for the two components.} \Cref{tab:2x2} crosses the memory with the readout on CIL-CORD. On every seed, the memory alone lifts full-label F1 over naive training by $\nAblMemMin$ to $\nAblMemMax$ points, and the blocked readout alone by \nAblReadMin{} to \nAblReadMax; together they lift it by $\nAblBothMin$ to $\nAblBothMax$ points, far more than the two add up to. This interaction is consistent with replay reducing drift and the readout reducing supervision conflict.

\paragraph{Ablation study for the background logit.} With every background logit fixed at zero, as in ICE-O, the memory and the blocked readout stay above CoLaR and below joint training on every seed and recover most of CoLaR-SBR's gain, while the learned background adds on two of three seeds (\cref{tab:ladder}). On BROS, the learned background is ahead on every seed ($\vBrosSbrFull$ against $\vBrosConstOFull$).

\paragraph{Ablation study for the memory (H4).} On CIL-CORD, SBR-Recall reaches $\vRecallFull$ full-label F1 ($\vRecall$ AA) with no stored document, within $\nRecallJointGap$ points of joint training on every seed, at the price of one inference pass per session. Our port of ICE-O (frozen encoder, constant Other logit) forgets nothing but reaches only $\vIceOFull$, below SBR-Recall on every seed. On Mixed without a memory, SBR-Recall scores $\nRecallMixNoMem$ AA, below naive fine-tuning (\cref{tab:main}): as predicted, the blocks of the other schemas fall silent until their domain returns, and with the $50$-document memory restored it scores $\nRecallMixMem$ on one seed (\cref{app:results}).

\paragraph{Memory analysis.}\label{sec:cost} In this section, we study how the memory is used (\cref{tab:memcmp,tab:memory}). First, a memory is needed: without it, our readout reaches only $\vNoBankFull$ full-label F1. Second, a memory alone is not enough: the same CoLaR memory under a shared softmax stays at $\vColarFull$, and raw replay of the same $250$ documents is no better (ER: $\vErFull$). Third, our readout is not limited to the latent memory: raw documents under the same readout reach $\vErSbrFull$ against $\vSbrFull$, with seed-matched differences of \nErSbrDiffFullMin{} to \nErSbrDiffFullMax{} points. Hence, the readout carries the gain, while the latent memory stores no pixels or token ids and trains only $\nMemTrainSbrPct\%$ of the parameters that raw replay trains. CoLaR keeps the $250$ documents in $\nMemColarMB$\,MB (\cref{app:memory}).
\section{Conclusion}
\label{sec:limitations}
\label{sec:conclusion}
In this paper, we addressed catastrophic forgetting in continual key-information extraction when the label space grows. We found that replay reduces the drift of the model, while a shared softmax still counts every background label against all classes, including those its annotator could not see. Based on this finding, we proposed a novel method called CoLaR-SBR, which maintains a compressed latent memory of whole documents to remind the model of past sessions and gives each session its own background logit to remove this conflict. We benchmarked our method against eight continual learning baselines on five public document datasets. The experiments show significant improvement in the class-incremental setting: on CIL-CORD, CoLaR-SBR achieves $\nSbrFull$ full-label F1, compared to $\nColarFull$ for the same memory with a shared softmax, and comes close to joint training ($\nJointFull$). It is noted that our study focuses on token classifiers with LayoutLMv3 as the main encoder, and that SBR-Recall adds one low-rank update and one inference pass per session. In the future, we wish to extend our method to large language models for document understanding, to other encoders and longer streams, and to other domains wherein the background label changes its meaning over time, such as incremental segmentation, object detection and named-entity recognition.\label{mainend}% last line of the main text (make_submission.sh checks its page)

\section*{LLM Usage Statement}
Claude (Anthropic), an LLM-based assistant, was used in this work to support running and monitoring experiments and to search and summarise related literature. All experiments were designed and decided by the authors, every reported number is produced by scripts from saved run files, and the authors reviewed and take full responsibility for all content.
{
    \small
    %% Overleaf free plan = 10 s compile limit. The bibliography is read from
    %% refs_precompiled.bbl whenever that file exists, so NO BibTeX pass runs on Overleaf even
    %% if refs.bib was uploaded too (BibTeX + its reruns alone exceed the limit). Both commands
    %% sit inside the conditional: a bare \bibliographystyle makes latexmk run BibTeX, which
    %% would overwrite the jobname .bbl with an empty one.
    %% Rule: the precompiled .bbl wins whenever it exists, so Overleaf runs
    %% pdflatex only (2 passes) even if refs.bib was uploaded too. To regenerate the
    %% bibliography locally: mv refs_precompiled.bbl aside, latexmk -pdf main.tex, then
    %%   cp main.bbl refs_precompiled.bbl   and re-upload main.tex + refs_precompiled.bbl.
    \IfFileExists{refs_precompiled.bbl}{\input{refs_precompiled.bbl}}%
      {\bibliographystyle{ieeenat_fullname}\bibliography{refs}}
}

%% Appendix heading without the paper title (one document: the body's \cref pointers resolve to it).
\ifwithappendix\clearpage\twocolumn[\centering\Large\textbf{Appendix}\vspace{1.0em}]\appendix\input{supp_body}\fi
\end{document}
